%% Согласно ГОСТ Р 7.0.11-2011:
%% 5.3.3 В заключении диссертации излагают итоги выполненного исследования, рекомендации, перспективы дальнейшей разработки темы.
%% 9.2.3 В заключении автореферата диссертации излагают итоги данного исследования, рекомендации и перспективы дальнейшей разработки темы.
Итогом выполнения диссертационной работы на соискание ученой степени кандидата технических наук, выполненной с целью совершенствования методики передачи координат в широком диапазоне расстояний при использовании современных методов обработки спутниковых данных, а также исследование различных факторов, влияющих на точность определения координат по РРР-алгоритму, является разработанная методика по созданию геодезических сетей при использовании РРР-алгоритма обработки спутниковых данных.

Основные научные и практические результаты заключаются в решении следующих актуальных научно-технических задач:

\begin{enumerate}
	\item Выполнено теоретическое и экспериментальное обоснование возможности использования РРР-алгоритма для построения геодезических сетей в различных диапазонах расстояний, которое позволило.......
	\item Проведены теоретические и экспериментальные исследования с целью выбора оптимальных параметров для обработки данных по РРР-алгоритму, которые позволили..........
	\item выполненные теоретические и экспериментальные исследования с целью выявления факторов, влияющих на точность обработки данных по РРР-алгоритму.
	
%	информационно-аналитический обзоров существующих методов создания геодезических сетей с использованием спутниковой аппаратуры. Приводятся основные требования по созданию геодезических сетей на примере ВГС (высокоточной геодезической сети) и СГС-1 (спутниковой геодезической сети).
%	\item Проведено научное исследование по выявлению влияния факторов на точность обработки данных по РРР-алгоритму. В качестве оцениваемых факторов были выделены следующие: продолжительность измерений; количество используемых спутников в процессе обработки по РРР-алгоритму; конфигурации различных ГНСС; влияние pdop фактора на точность; влияние дискретности (частоты записи).
%	\item Разработана методика создания геодезических сетей с использованием РРР-алгоритма обработки спутниковых данных. Методику можно применять при построении геодезических сетей, в том числе государственной геодезической сети. В ходе экспериментальных исследований было установлено, что точность определения координат с использованием РРР-алгоритма получается несколько выше, чем при классической обработке в статике.
\end{enumerate}
%
%	\item Теоретическое и экспериментальное обоснование возможности использования РРР-алгоритма для построения геодезических сетей в различных диапазонах расстояний.
%	\item Теоретические и экспериментальные исследования с целью выбора оптимальных параметров для обработки данных по РРР-алгоритму.
%	\item Теоретические и экспериментальные исследования с целью выявления факторов, влияющих на точность обработки данных по РРР-алгоритму.

Результаты диссертационного исследования рекомендуются к использованию при поддержании и развитии ГГС, а также для дальнейшего исследования точностных возможностей определения координат при использовании РРР-алгоритма.

Перспективы дальнейших исследований по данному направлению заключаются .......



